\chapter{The Interval Renderer (Main Pathway)}\label{ch:interval-renderer}

\section{Overview}

The interval shader is the production renderer used by the interactive viewer
(default \code{RenderMode::IntervalShader}) and by forward-only inference.
Its central trick: instead of asking the fragment shader to figure out where
a ray enters and exits a tetrahedron (which requires per-fragment plane
intersections, as the legacy renderers do), a compute pass \emph{decomposes
each tet's screen projection into non-overlapping triangles} whose vertices
carry the front and back NDC depths of the tet along that pixel's ray. The
hardware rasterizer then interpolates the depth interval per pixel, and each
covered pixel is shaded exactly once per tet with the correct entry/exit
segment. Overdraw is resolved by ordinary premultiplied-alpha blending over
back-to-front sorted tets --- no fragment sorting, no fixed-function tricks.

Code lives in \code{crates/rmesh-render} (\code{src/lib.rs} and
\code{src/wgsl/}). Two pipeline structs drive the path:

\begin{itemize}
\item \code{ComputeIntervalPipelines} --- compute generation followed by a
  vertex/fragment draw. \textbf{Universal}: works on Metal, Vulkan, DX12 and
  WebGPU. This is the production path.
\item \code{IntervalPipelines} --- a single mesh-shader dispatch
  (\code{interval_mesh.wgsl}) replacing the generation + vertex stages.
  Requires \code{Features::EXPERIMENTAL_MESH_SHADER}; currently untested for
  lack of hardware.
\end{itemize}

\section{Stage 1 --- Projection, culling, SH color}

\code{project_compute.wgsl} (or the lean \code{project_compute_hw.wgsl} when
tile counts are not needed) runs one thread per tet
(\code{@workgroup_size(64)}, 2D-dispatched above 65\,535 workgroups):

\begin{enumerate}
\item Project the 4 vertices to NDC. If any \code{clip.w <= 0} the tet
  straddles the camera plane; it is kept and near-clipping is deferred to the
  rasterizer. Otherwise reject if the NDC AABB is fully outside
  $[-1,1]^3$ or the screen extent is under one pixel${}^2$.
\item Cull by maximum optical depth: skip if
  $\rho \cdot \max_{\text{edge}} < 0.01$.
\item Compute the depth sort key from the circumsphere power
  (see Chapter~\ref{ch:gpu-sorting}), bit-inverted so ascending order is
  back-to-front. Culled threads write \code{0xFFFFFFFF}, sorting to the end.
\item An \code{atomicAdd} on the instance count in \code{indirect_args}
  compacts the visible tets into \code{compact_tet_ids}.
\item Evaluate SH color (degree 0--3, f16 coefficients unpacked with
  \code{unpack2x16float}) in the direction
  $\widehat{\text{centroid} - \text{cam}}$, add the $+0.5$ DC bias and a
  gradient offset, then \code{softplus}; write to \code{colors[tet_id]}.
  Without SH data (training), base colors are copied through.
\end{enumerate}

\code{project_compute_hw.wgsl} additionally binds \code{Uniforms} as a
\emph{uniform} buffer (dropping one storage binding for WebGPU's cap) and
skips the tile-touch counting loop, which dominates the cost of the full
variant. \code{project_compute_16bit.wgsl} quantizes depth to 16 bits for the
faster \code{SortMode::Gpu16} (4 radix passes instead of 8).

\section{Stage 2 --- Radix sort}

\code{rmesh-sort} sorts \code{(sort_keys, sort_values)} by depth. The DRS
subgroup backend is used when \code{Features::SUBGROUP} is available,
otherwise the Basic 5-pass backend; on backends without either (Chrome
WebGPU) or with \code{SortMode::Cpu}, a CPU sorter does frustum culling and
power sorting and its output is uploaded with a \code{StagingBelt} between
the projection and generation submits. Details in Chapter~\ref{ch:gpu-sorting}.
\code{record_radix_sort} returns \code{result_in_b}, which selects the A/B
bind group downstream (the sort ping-pongs buffers).

\section{Stage 3 --- Indirect convert}

\code{interval_indirect_convert.wgsl} runs a single thread that reads the
visible count and fills an 8$\times$u32 buffer \code{interval_args_buf}:
words $[0..2]$ are compute dispatch args
$(\lceil n/64 \rceil, 1, 1)$ for the generation pass, words $[3..7]$ are
\code{draw_indexed_indirect} args $(12n, 1, 0, 0, 0)$ --- twelve indices per
tet. No CPU readback of the count is ever needed.

\section{Stage 4 --- Interval decomposition (\code{interval_compute.wgsl})}

One thread per sorted visible tet. The screen projection of a tetrahedron has
exactly two silhouette topologies, and both are handled with a small,
fixed-size output:

\paragraph{Case 1 --- one vertex projects inside the opposite face.}
\code{classify_silhouette} tests each vertex against the opposite face's
screen triangle with signed 2D cross products. The silhouette is that face's
triangle; the interior vertex splits it into a 3-triangle fan around the
\emph{center} point. At the center, the ray passes through both the face
(depth $z_{\text{face}}$, barycentric-interpolated) and the interior vertex
(depth $z_{\text{center}}$); $z_{\text{front}} = \min$,
$z_{\text{back}} = \max$. Emits 4 vertices, 3 triangles.

\paragraph{Case 2 --- no interior vertex.} Two opposite edges cross in
screen space. \code{find_crossing_edges} tests the three opposite-edge pairs
$(01,23), (02,13), (03,12)$ with a 2D line--line intersection requiring both
parameters in $(0,1)$. The silhouette is a quad of the 4 vertices, split into
a 4-triangle fan around the crossing point; the two crossing edges provide
the front and back depths there. Emits 5 vertices, 4 triangles.

Per-vertex color offsets $\nabla c \cdot (v - o_{\text{cam}})$ are carried
separately for the front and back surfaces, and the scalar-field gradient
(used later as a normal) is carried for both. Degenerate tets (behind camera,
depth range $< 10^{-7}$, density $< 10^{-6}$) write zeroed vertices.

\paragraph{Fixed-slot output.} Every tet owns exactly 5 vertex slots
(4 \code{vec4} each, 320 bytes/tet) in \code{interval_vertex_buf} at base
$5\,\code{tid}$, plus two \code{vec4} of flat data
(\code{base_color}, \code{density}, \code{tet_id}) in
\code{interval_tet_data_buf}:
\begin{lstlisting}
slot + 0 : (ndc_x, ndc_y, z_front, z_back)
slot + 1 : (color_off_front, color_off_back, 0, 0)
slot + 2 : (gradient_front.xyz, 0)
slot + 3 : (gradient_back.xyz, 0)
\end{lstlisting}
Because the layout is fixed, no atomic compaction is needed: a \emph{static
fan index buffer} (\code{interval_fan_index_buf}) built once at load indexes
the slots directly --- per tet $i$ with $b = 5i$, the triangles
\[
(b,\,b{+}1,\,b{+}4),\quad (b{+}1,\,b{+}2,\,b{+}4),\quad
(b{+}2,\,b{+}3,\,b{+}4),\quad (b{+}3,\,b,\,b{+}4).
\] Case-1 tets duplicate slot~0 into slot~3, making
the fourth fan triangle degenerate.

\section{Stage 5 --- Vertex shader (\code{interval_vertex.wgsl})}

The pipeline binds \emph{no} vertex buffers; the vertex shader derives
\code{tet = vertex_index / 5} and reads the slots from storage. It forwards
NDC position, the depth pair, the color-offset pair, flat density /
base-color / tet-id, and the interpolated field gradient.

\section{Stage 6 --- Fragment shader (\code{interval_fragment.wgsl})}

Per pixel:
\begin{enumerate}
\item Linearize the interpolated NDC depths:
  $z = \dfrac{n f}{f - z_{\text{ndc}}(f - n)}$.
\item Reconstruct the ray scale from the intrinsics,
  $s = \lVert(\frac{x-c_x}{f_x}, \frac{y-c_y}{f_y}, 1)\rVert$, giving the
  chord length $d = |z_b - z_f|\, s$.
\item Entry/exit colors $c_{f/b} = \max(c_{\text{base}} + \text{off}_{f/b}, 0)$.
\item Apply the closed-form volume integral (see the overview article):
  $\tau = \rho d$, $\alpha = 1 - e^{-\tau}$,
  $c_{\text{premul}} = w_0 c_b + w_1 c_f$.
\end{enumerate}

Two entry points share the shading core: \code{main_color_only} writes just
the color target (lean forward path), and \code{main} writes the full MRT
G-buffer for deferred shading (all channels premultiplied by $\alpha$ for
hardware blending):

\begin{longtable}{@{}lll@{}}
\toprule
Location & Format & Contents \\
\midrule
0 & Rgba16Float & $\text{albedo}\cdot\alpha,\ \alpha$ \\
1 & Rgba16Float & $\text{roughness}\cdot\alpha,\ f_{0}\cdot\alpha,\ \text{metallic}\cdot\alpha,\ \alpha$ \\
2 & Rgba8Unorm & bias-encoded field-gradient normal $\cdot\,\alpha$ \\
3 & Rgba16Float & $\alpha E[z],\ 0,\ \alpha E[z^2],\ \alpha$ (expected termination depth) \\
\bottomrule
\end{longtable}

$E[z^2]$ is exact for the exponential termination distribution within the
segment, with a Taylor correction
$\tau/6 - \tau^2/12 + \tau^3/40$ to avoid cancellation. The blend state is
premultiplied over ($\text{src}\cdot 1 + \text{dst}\cdot(1-\alpha_\text{src})$),
depth test \code{LessEqual} against the primitive depth with depth writes
disabled, and no face culling (the triangles are screen-aligned).

\section{The mesh-shader variant}

\code{interval_mesh.wgsl} fuses stages 4--5 into one mesh-shader dispatch:
\code{@workgroup_size(16)}, one thread per tet, 16 tets per workgroup. Each
thread performs the same silhouette classification, then a serial 16-element
workgroup prefix sum compacts variable per-tet output (4~verts/3~prims or
5~verts/4~prims) into the workgroup's mesh output arrays
(\code{MAX_VERTS=240}, \code{MAX_PRIMS=192}). A hard-won implementation note
from the source: divergent threads must \emph{never} early-return in a mesh
shader (it hangs); output is guarded with \code{if my_vert_count > 0u}
instead. Dispatch args come from \code{indirect_convert.wgsl} with
\code{TETS_PER_GROUP=16} and \code{draw_mesh_tasks_indirect}.

\section{Record functions and orchestration}

In \code{crates/rmesh-render/src/lib.rs}:
\begin{itemize}
\item \code{record_sorted_compute_interval_forward_pass}\\
  --- the GPU-sort
  path: reset args $\to$ project $\to$ radix sort $\to$ indirect convert
  $\to$ generation compute (\code{dispatch_workgroups_indirect}) $\to$
  MRT/color render pass (\code{draw_indexed_indirect} at byte offset 12).
\item \code{record_compute_interval_project} /
  \code{record_compute_interval_after_sort} --- the split used by the CPU
  sort path, with a \code{StagingBelt} write of sorted indices and the
  visible count in between.
\item \code{record_sorted_interval_forward_pass} --- the mesh-shader
  variant.
\end{itemize}

The viewer (\code{crates/rmesh-viewer/src/render.rs}) selects the pass by
\code{RenderMode} and \code{SortMode}; all volume passes use
\code{LoadOp::Load} to composite over the earlier opaque-primitive pass.
The same generation + moment-render machinery is reused from each light's
perspective to build deep shadow maps (Chapter~\ref{ch:shadow-rendering}).

